% Results section.
%
% Every number in this file that is not in a table is quoted from one, and
% every table under tables/ is generated by cf_worldmodels/experiments/
% export_tables.py from the stored metrics.json files. Regenerate with:
%
%     cd cf_worldmodels && python experiments/export_tables.py
%
% Requires booktabs.

\section{Results}
\label{sec:results}

We ran the full grid: three environment families $\times$ three dynamic-distance
levels $\times$ five methods $\times$ five seeds, for 225 runs, each training a
world model on task~A and then on task~B under one protocol (5000 gradient
steps, 20 random-policy episodes per task, batch~8, sequence length~5). A
further 45 reference pairs --- one model per environment, trained from scratch
--- supply the from-scratch arm of forward transfer and the measured distance
$d_{\mathrm{trans}}$. No run dropped a NaN step.

\subsection{How the numbers are reported}
\label{sec:results:reporting}

Cells are summarised by the \emph{median} over their five seeds, with the full
seed range where space allows. This is not cosmetic. RD is a KL divergence and
is unbounded above, so a method whose post-task-B model becomes overconfident
acquires a heavy right tail by construction: our own method, UG-MTM on
MiniGrid at maximum distance, produces $17.7$, $40.0$, $520$, $574$ and $4364$
across its five seeds, and the mean of those, $1103$, describes none of them.
The skew is common rather than exceptional --- 22 of the 180 forgetting cells
have an upper half that stretches more than five times further from the median
than their lower half --- so we report medians throughout and mark the skewed
cells ($\ast$) rather than averaging the tail away.

PF and RD are reported separately. The aggregate WMF of the earlier
formulation is dominated by RD, which supplies 78--97\% of it, so a single
scalar would be RD with extra steps. Comparisons between methods are paired by
seed and tested with an exact two-sided permutation test over all $2^5$ sign
flips. With $n=5$ that test cannot return a $p$ below $2/2^5 = 0.0625$; we
report it alongside the paired effect size $d_z$ and the win count, and we do
not report parametric $p$-values, which at this $n$ would be statements about
the normality assumption rather than about the data.

\subsection{The labelled distance axis does not order forgetting}
\label{sec:results:axis}

The benchmark was built around an ordinal axis: each family has a minimum, a
medium and a maximum dynamic distance, and the expectation encoded in that
design is that forgetting grows along it. It does not.

\begin{table}[t]
  \centering
  % Generated by experiments/export_tables.py -- do not edit by hand.
% Regenerate after any change to results/ or to the aggregation.

\begin{tabular}{lrrrc}
\toprule
Family & min & med & max & Peak \\
\midrule
MiniGrid & 33.31 & \textbf{62.26} & 35.63 & med \\
Gymnasium & 77.45 & \textbf{85.91} & 59.59 & med \\
DMControl & 1.15 & \textbf{70.94} & 52.66 & med \\
\bottomrule
\end{tabular}

  \caption{Rollout divergence by family and labelled distance level, aggregated
  over the four methods that share the RSSM architecture (median over methods
  of each method's median over seeds). Forgetting peaks at the \emph{medium}
  level in all three families.}
  \label{tab:axis}
\end{table}

Table~\ref{tab:axis} gives RD per family and level. In all three families the
maximum sits at the \emph{medium} level, and the labelled order
$\text{min} < \text{med} < \text{max}$ is therefore wrong in every one of them.
Three out of three is the most robust part of this result: it does not depend on
how cells are aggregated, on which of the two forgetting metrics is read, or on
any threshold.

Gymnasium makes the point without any appeal to how the levels compare across
families, because there the two upper levels are nested by construction. Both
move HalfCheetah from gravity $9.8$ to $4.0$; the maximum level additionally
scales mass by $3$ and friction by $0.5$. It is strictly the medium
perturbation plus more perturbation, and it produces \emph{less} forgetting
($59.59$ against $85.91$).

The obvious objection is that this is a statement about our budget rather than
about forgetting: if at the maximum level the model never fits task~B, nothing
gets overwritten and the cell reads low for an uninteresting reason. We tested
it by rerunning both cells at twice the training budget, and it does not hold
up. The ordering survives --- restricted to fine-tuning, so that the two budgets
are compared on the same estimator, the cells read $118.04$ (medium) against
$58.79$ (maximum) at 5000 updates and $144.70$ against $102.18$ at 10000 --- and
the mechanism the objection proposes is the opposite of what happens. Task~B is
\emph{easier} to fit at the maximum level than at the medium one, at both
budgets ($17.31$ against $27.10$, and $15.14$ against $24.65$): the heavier,
lower-friction, low-gravity cheetah moves less, and a less varied task is a
task with less to learn from. What we report as robust is the direction, which
survives the doubling; the magnitude is budget-dependent, the ratio between the
two cells falling from $2.01$ to $1.42$.

The aggregate in Table~\ref{tab:axis} deliberately excludes UG-MTM. It freezes
its encoder, which puts its RD two orders of magnitude below the other four
methods in DMControl ($0.002$ against $1.2$); pooling all five would average two
scales rather than read the axis. Its own numbers appear in
Section~\ref{sec:results:methods} and in Table~\ref{tab:rd}.

\subsection{What does predict forgetting}
\label{sec:results:predictors}

If the labelled level does not order forgetting, something else may. We compare
three candidate predictors across the nine cells: the labelled level itself, the
measured distance $d_{\mathrm{trans}}$ between the two environments' transition
models, and the difficulty of task~B --- operationalised as the held-out
reconstruction error a model reaches on task~B after training on it, i.e.\ how
hard task~B is to fit at this budget.

\begin{table}[t]
  \centering
  % Generated by experiments/export_tables.py -- do not edit by hand.
% Regenerate after any change to results/ or to the aggregation.

\begin{tabular}{llrrr}
\toprule
Family & Level & RD & $d_{\mathrm{trans}}$ & Task-B difficulty \\
\midrule
MiniGrid & min & 33.73 & 22.35 & 2.01 \\
 & med & 60.88 & 20.64 & 42.60 \\
 & max & 32.11 & 25.19 & 23.59 \\
\addlinespace
Gymnasium & min & 77.45 & 28.75 & 24.63 \\
 & med & 85.91 & 30.86 & 27.28 \\
 & max & 64.01 & 24.08 & 17.44 \\
\addlinespace
DMControl & min & 1.15 & 1.87 & 15.43 \\
 & med & 71.64 & 7.58 & 34.07 \\
 & max & 52.61 & 7.46 & 40.54 \\
\midrule
\multicolumn{2}{l}{Spearman $\rho$ with RD} & 0.00 & 0.58 & 0.43 \\
\multicolumn{2}{l}{\quad (predictor)} & \emph{labelled level} & \emph{measured} & \emph{measured} \\
\bottomrule
\end{tabular}

  \caption{The nine cells behind the rank correlations, and the correlations
  themselves. RD and task-B difficulty are aggregated as in
  Table~\ref{tab:axis}; $d_{\mathrm{trans}}$ is a property of the task pair and
  is the median over seeds.}
  \label{tab:predictors}
\end{table}

Table~\ref{tab:predictors} shows the outcome. The labelled level is
uninformative ($\rho = 0.05$). The measured distance does better
($\rho = 0.53$), and the difficulty of task~B does best ($\rho = 0.58$). In
MiniGrid and Gymnasium the two columns rise and fall together, and the medium
level is exactly where task~B is hardest to fit ($42.26$ against $2.35$ and
$21.53$ in MiniGrid).

The reading we draw is that \textbf{forgetting scales with how much the new task
demands of the model, not with how far it is from the old one}. We state it as
the best available account rather than a demonstrated mechanism, because
DMControl does not follow it: task-B difficulty there keeps rising from minimum
to maximum ($15.43$, $35.17$, $39.00$) while RD falls after the medium level.
Task-B difficulty is the best of the three predictors and it explains two
families out of three; what all nine cells agree on is the negative result, that
the labelled axis explains none.

Two caveats belong with this. First, $n = 9$, with three families on very
different scales, and a rank correlation over pooled cells is indicative rather
than conclusive; the repeated peak of Section~\ref{sec:results:axis} is the
robust half of the finding and the correlations are the interpretation of it.
Second, and more seriously, $d_{\mathrm{trans}}$ is not a drop-in replacement
for the labelled axis, and its $\rho = 0.53$ should not be read as one. Between
families it separates cleanly --- DMControl's cells sit near $5$ where
MiniGrid's sit near $20$ --- but \emph{within} a family it mostly does not.
Gymnasium's three levels have medians of $28.75$, $30.86$ and $24.71$ spanning
about six units, while each of those cells spans fifteen to nineteen across its
five seeds ($[21.5, 37.0]$, $[18.5, 37.2]$, $[17.7, 30.1]$); the rank order is
a rank order over medians that are inside each other's noise. MiniGrid inverts
min and med on the same terms ($20.07$ against $19.59$). The budget rerun of
Section~\ref{sec:results:axis} confirms the diagnosis directly: at twice the
updates, Gymnasium's medium and maximum medians swap places ($57.30$ against
$71.51$), while their seed ranges remain all but identical. Much of the rank
correlation is therefore carried by between-family separation, which is the
comparison $d_{\mathrm{trans}}$ is least entitled to make
(Section~\ref{sec:method:distances}).

What survives is narrower than a replacement axis and still worth stating: a
measured distance is answerable to evidence in a way a label is not, and this
one answered. It says its own within-family orderings are not resolved at five
seeds, and it changes when the budget changes --- which is itself a finding
about Equation~\ref{eq:dtrans}, since a distance between two environments
should not depend on how long one trains the models used to measure it.

\subsection{Three cells produce no forgetting, and we report them as controls}
\label{sec:results:controls}

Not every cell has something to forget. We call a cell a \emph{control} when no
method's held-out task-A reconstruction error grows by more than 10\% after
training on task~B. The threshold is not a knife edge: in every cell that does
produce forgetting, the worst method ends between $15\times$ and $753\times$ its
original error, so the nearest one sits more than $13\times$ above the threshold
and any cutoff between $1.1$ and $14$ selects the same three cells.

Three of the nine qualify (marked $\dagger$ in Table~\ref{tab:encoder}):
Gymnasium at minimum and medium distance, and DMControl at minimum. In the first
two the task pair is the same HalfCheetah under different gravities, so training
on task~B keeps teaching the model to reconstruct task~A. This is what the bottom
of a distance axis is supposed to look like, and the benchmark detects it; it
also means those cells discriminate between methods on nothing, and the effective
grid is six cells, not nine.

The criterion is deliberately blind to RD, and Gymnasium's medium cell shows
why: it loses nothing in pixels and simultaneously carries the highest RD in its
family ($85.91$). A cell can leave the encoder intact and still move the
transition component a long way. That is the subject of the next section.

\subsection{Forgetting happens in the encoder, where the metrics cannot see it}
\label{sec:results:encoder}

\begin{table}[t]
  \centering
  % Generated by experiments/export_tables.py -- do not edit by hand.
% Regenerate after any change to results/ or to the aggregation.

\begin{tabular}{llrrrrr}
\toprule
Family & Level & Fine-tuning & Replay & EWC & Prog.\ Nets & UG-MTM \\
\midrule
MiniGrid & min & 811 & 1.31 & 800 & 797 & 1.00 \\
 & med & 829 & 1.18 & 848 & 825 & 1.00 \\
 & max & 32 & 1.02 & 31 & 31 & 1.00 \\
\addlinespace
Gymnasium & min$^{\dagger}$ & 0.94 & 0.90 & 0.94 & 0.97 & 1.00 \\
 & med$^{\dagger}$ & 1.04 & 0.92 & 1.05 & 1.06 & 1.00 \\
 & max & 14 & 1.04 & 15 & 14 & 1.00 \\
\addlinespace
DMControl & min$^{\dagger}$ & 1.02 & 1.00 & 1.02 & 1.02 & 1.00 \\
 & med & 16 & 0.37 & 16 & 16 & 1.00 \\
 & max & 17 & 0.51 & 17 & 16 & 1.00 \\
\bottomrule
\end{tabular}

  \caption{Held-out task-A reconstruction error after training on task~B, as a
  factor of the same model's error before it. $1.00$ means nothing was lost,
  $753$ means the error grew by a factor of 753. $\dagger$ marks control cells
  (Section~\ref{sec:results:controls}).}
  \label{tab:encoder}
\end{table}

PF and RD are computed on latents encoded once, by the post-task-A model, which
is what isolates the transition component $M$ --- the benchmark's declared
scope. The cost of that isolation is that both metrics are blind to drift in the
encoder itself, and Table~\ref{tab:encoder} shows how large the blind spot is.

Fine-tuning's held-out task-A reconstruction grows by a factor of $745$ in
MiniGrid at minimum distance, and its PF for that cell is $-6.21$: measured in
the frozen latent basis, the model \emph{improved}. This is not an isolated
sign error. Fine-tuning's PF is negative in eight of the nine cells
(Table~\ref{tab:pf}) while its encoder is destroyed in six of them.

EWC makes the dissociation sharpest, and it was predicted before it was
measured. Its Fisher information is defined over $\log P(z' \mid z, a)$, so it
is \emph{exactly} zero on every encoder parameter: the penalty cannot constrain
the VAE even in principle. Both halves of that prediction hold. In MiniGrid at
minimum distance EWC preserves the transition component almost perfectly
($\text{PF} = -0.07$, against $-6.21$ for fine-tuning) while its encoder
degrades by a factor of $753$ --- marginally worse than fine-tuning's $745$.
Across the grid its PF stays within $|0.07|$ in seven of the nine cells.

Replay is the only method that keeps both. Its task-A reconstruction never grows
by more than $27\%$, and in DMControl it \emph{improves} (factors of $0.37$ and
$0.45$), which is what training on A and B together should do.

We report this as a scope declaration rather than a defect: measuring $M$ in a
fixed latent basis is the stated purpose, and the alternative --- letting each
model encode with its own drifting VAE --- would compare each model against its
own moving coordinates rather than against a common evaluation set. But a
benchmark that reports only PF and RD would report that fine-tuning does not
forget, and that is why every run in this suite records the pixel scale
alongside them.

\subsection{What the benchmark says about the methods}
\label{sec:results:methods}

\begin{table}[t]
  \centering
  % Generated by experiments/export_tables.py -- do not edit by hand.
% Regenerate after any change to results/ or to the aggregation.

\begin{tabular}{llrrrrr}
\toprule
Family & Level & Fine-tuning & Replay & EWC & Prog.\ Nets & UG-MTM \\
\midrule
MiniGrid & min & -5.58 & -5.98 & -0.03 & 4.06 & 3.25 \\
 & med & -4.23 & -7.23 & 0.04 & 6.10 & 11.71 \\
 & max & 0.60 & -6.32 & 0.07 & 14.43 & 38.65 \\
\addlinespace
Gymnasium & min & -8.27 & -8.00 & 0.02 & -0.03 & -0.03 \\
 & med & -8.40 & -8.27 & -0.01 & -0.36$^{\ast}$ & -0.05 \\
 & max & 0.52 & -6.00 & 0.04 & 15.09 & 15.66 \\
\addlinespace
DMControl & min & -5.43 & -5.23 & 0.01 & 5.80 & 0.11 \\
 & med & -18.02 & -13.54 & -11.11$^{\ast}$ & -3.05 & 0.47 \\
 & max & -16.95 & -13.39 & -9.56 & -3.53 & 2.08 \\
\bottomrule
\end{tabular}

  \caption{PF (prediction fidelity) by method and cell; median over five seeds.
  Positive means the transition component got worse at task~A. $\ast$ marks
  cells whose seed sample is right-skewed.}
  \label{tab:pf}
\end{table}

\begin{table}[t]
  \centering
  % Generated by experiments/export_tables.py -- do not edit by hand.
% Regenerate after any change to results/ or to the aggregation.

\begin{tabular}{llrrrrr}
\toprule
Family & Level & Fine-tuning & Replay & EWC & Prog.\ Nets & UG-MTM \\
\midrule
MiniGrid & min & 41.35 & 24.94 & 28.42 & 38.19 & 12.90 \\
 & med & 72.00$^{\ast}$ & 37.43 & 59.52 & 65.01$^{\ast}$ & 32.33 \\
 & max & 47.28 & 23.41 & 23.99 & 51.47 & 520$^{\ast}$ \\
\addlinespace
Gymnasium & min & 75.30 & 82.62 & 39.42 & 79.59 & 10.17 \\
 & med & 118 & 78.32 & 74.45 & 93.50 & 28.26 \\
 & max & 58.79 & 90.61 & 42.21 & 60.40 & 139$^{\ast}$ \\
\addlinespace
DMControl & min & 1.19 & 1.11 & 0.00$^{\ast}$ & 1.54$^{\ast}$ & 0.00 \\
 & med & 78.78 & 66.37 & 71.74 & 70.15 & 0.01 \\
 & max & 59.44 & 47.62 & 57.71 & 26.85 & 0.14 \\
\bottomrule
\end{tabular}

  \caption{RD (rollout divergence) by method and cell; median over five seeds.
  $\ast$ marks right-skewed cells.}
  \label{tab:rd}
\end{table}

\paragraph{Replay is not insufficient here.} On the six cells that produce
forgetting, infinite replay has lower RD than fine-tuning in five, unanimously
across seeds in four of them ($p = 0.0625$, the floor at $n=5$; $d_z$ from
$-1.99$ to $-3.05$). The exception is Gymnasium at maximum distance, where
replay is \emph{worse} in all five seeds ($90.61$ against $58.79$,
$d_z = +1.79$). The direction we measure is the one expected a priori, since
replay trains on A and B together; we note it because an earlier version of this
benchmark reported the opposite, at $p < 0.001$ with five seeds, and that
reversal is attributable to defects in the measurement rather than to the
setting.

\paragraph{EWC protects the transition component and only that.} Discussed in
Section~\ref{sec:results:encoder}. Its near-zero PF is not universal: in
DMControl's two forgetting cells it reads $-11.32$ and $-9.27$, in line with the
other methods, so the protection is not uniform across families either.

\paragraph{Progressive networks are the only baseline with consistently
positive PF} in MiniGrid ($3.08$, $5.97$, $13.50$), i.e.\ the benchmark scores
them as forgetting more of the transition component, not less --- while their
encoder degrades like fine-tuning's. Column capacity does not help a component
that is not the one being overwritten.

\paragraph{UG-MTM does not forget because it does not learn.} Its encoder factor
is exactly $1.00$ in all nine cells --- it is frozen, so task-A reconstruction is
preserved bit for bit --- and it pays for that in forward transfer: FT of $-559$,
$-525$ and $-1603$ in MiniGrid (Table~\ref{tab:ft}), where no other method
leaves the range $[-10.48, +3.40]$ anywhere in the grid. The cost nearly
vanishes when the two tasks
look alike ($-4.33$ in Gymnasium at minimum distance), which locates the regime
where structural isolation is affordable rather than establishing that it works.

\paragraph{Freezing the encoder produces overconfident transition models.} This
is the benchmark's own finding about our method, and it is why
Section~\ref{sec:results:reporting} reports medians. UG-MTM's RD in MiniGrid at
maximum distance spans $17.7$ to $4364$ over five seeds. Reproducing the extreme
seed and instrumenting the KL term by term shows the cause is not a diverging
rollout: the KL is already $80$ at rollout step~0, the predicted means barely
move over fifteen steps, and what explodes is the post-task-B variance, with
$\sigma = 0.007$ in some latent dimensions. The term
$(\mu_i - \mu_k)^2 / \sigma_k^2$ does the rest. Confident and wrong is a failure
mode a rollout-divergence metric will expose whenever it is unbounded above, and
truncating the horizon would not hide it.

\begin{table}[t]
  \centering
  % Generated by experiments/export_tables.py -- do not edit by hand.
% Regenerate after any change to results/ or to the aggregation.

\begin{tabular}{llrrrrr}
\toprule
Family & Level & Fine-tuning & Replay & EWC & Prog.\ Nets & UG-MTM \\
\midrule
MiniGrid & min & -1.09 & -2.27 & -1.70$^{\ast}$ & -1.04 & -559 \\
 & med & -1.87 & -0.82 & -2.71 & -2.98 & -525 \\
 & max & 3.40 & -2.04 & 0.50 & -0.34 & -1603 \\
\addlinespace
Gymnasium & min & 1.83 & 0.33 & 1.94 & 1.80 & -4.33 \\
 & med & 1.92 & 0.12 & 2.01 & 1.81 & -14.49 \\
 & max & -3.24$^{\ast}$ & -10.48 & -2.67 & -2.88 & -251 \\
\addlinespace
DMControl & min & 0.12 & 0.05 & 0.12 & 0.19$^{\ast}$ & -0.33 \\
 & med & -2.29 & -7.84 & -1.91 & -1.04 & -285 \\
 & max & -3.09 & -8.54$^{\ast}$ & -3.62 & -1.77 & -314 \\
\bottomrule
\end{tabular}

  \caption{FT (forward transfer) by method and cell, in pixel space: held-out
  task-B reconstruction of a model trained on B from scratch minus that of the
  pretrained model. Positive means knowledge of task~A helped.}
  \label{tab:ft}
\end{table}

\subsection{Threats to validity}
\label{sec:results:threats}

\textbf{Five seeds.} An exact paired permutation test cannot return $p < 0.0625$
at $n=5$, so no comparison here reaches conventional significance without
leaning entirely on a normality assumption. We report effect sizes and win
counts for that reason.

\textbf{Nine cells.} The rank correlations of
Section~\ref{sec:results:predictors} pool cells from families whose scales
differ by an order of magnitude. They are the interpretation of
Table~\ref{tab:predictors}, not an independent test of it.

\textbf{Six discriminating cells.} Three of the nine are controls
(Section~\ref{sec:results:controls}).

\textbf{Budget.} 20 random-policy episodes per task and 5000 gradient updates.
Task~A is learned at this scale, which a forgetting benchmark has to show before
it can claim anything was forgotten: a scaling study on MiniGrid takes held-out
task-A reconstruction from $6.49$ at 1000 updates to $1.61$ at the 5000 used
here and $1.37$ at 10000, so the chosen budget is past the steep part of the
curve but not at its floor, and absolute magnitudes should be read as
budget-dependent. Two of the nine cells were rerun at twice the budget, and what
that establishes is bounded: the direction of the axis result survives in the
family where the levels are nested, and $d_{\mathrm{trans}}$ does not
(Sections~\ref{sec:results:axis} and~\ref{sec:results:predictors}). The other
seven cells were not rerun. The protocol is declared in the config, recorded in
every result file, and refused by the runner if two budgets would be averaged
into one cell.

\textbf{One task switch.} The formalism is stated for $T_1 \ldots T_k$; the
experiments run $k = 2$.
